\documentclass[10pt,twocolumn,letterpaper]{article}
\usepackage[margin=0.72in,columnsep=0.26in]{geometry}
\usepackage[T1]{fontenc}
\usepackage{times,microtype,amsmath,amssymb,booktabs,graphicx,array}
\usepackage[breaklinks,colorlinks,urlcolor=blue,citecolor=black,linkcolor=black]{hyperref}
\usepackage{url,balance}
\raggedbottom
\setlength{\parindent}{1em}
\setlength{\parskip}{0pt}
\newcommand{\art}[2]{\href{https://github.com/rz4/reticuli/blob/5318acaac34ef99eeba9ff064a995a6313cc70ec/#1}{#2}}
\newtheorem{condition}{Condition}
\newtheorem{claimx}{Claim}
\title{What Must Be in the Hash:\\Soundness Conditions for Content-Addressed\\Executable Specifications}
\author{Rafael Zamora-Resendiz}
\date{Working manuscript --- October 2026}
\begin{document}
\maketitle

\begin{abstract}
Name a piece of software by a hash over its acceptance criteria, excluding its implementation, and any implementation that passes the criteria carries the name. The construction is simple. Deciding \emph{which bytes belong in the hash} is not.

We give two conditions on the committed byte set. They press from opposite sides. Commit less than the acceptance procedure can read, and one identity admits contradictory verdicts. Commit more than the specification means, and implementations that satisfy every requirement are refused over presentation. Both failures occurred in a working system. We exhibit both: a gate that read producer hints its own hash omitted, and two conforming packagers that emitted the same steps in different orders and minted different names. The first repair is the instructive one. The read set shrinks; the hash does not grow, because growing it would violate the second condition.

We then extend identity to composed specifications, and state the rule a composed verdict must satisfy. Three working forgeries show what its clauses buy. A forged dependency attribution passed every verification verb in every implementation we tested. A signature validated a chain that the signer's own deep audit rejects. And a user-facing verdict silently resolved to the primitive that does not compose.

Two further results. A finite set of examples cannot constrain behavior outside its support: the class it admits contains a lookup table, and a blind reconstruction built exactly that. Obligations evaluated over inputs generated at judgment time are therefore necessary, not merely useful. And complexity obligations can sit inside a portable identity after all --- not as absolute bounds, which would make identity host-relative, but as growth ratios. One such obligation caught an exponential traversal that every functional criterion accepted.

The conditions are validated at their limit by a system that reconstructs itself: across three generations and two independent producers, one identity survived complete replacement of the implementation it names.
\end{abstract}

\section{Introduction}
Content addressing names data by a hash of its bytes. Naming \emph{software} this way is unsatisfying, because the thing one wants to name is rarely the bytes. Two implementations of the same requirements deserve the same name; the same implementation reformatted does not deserve a new one.

An alternative is to hash the \emph{requirements}. Let the acceptance criteria --- executable tests, declared interfaces, fixtures, pinned verdicts --- be the committed object, and let the implementation lie outside the commitment. The name then denotes an equivalence class: every implementation the criteria accept carries it. This is attractive for replacement, archival, and provenance, and the construction takes one line of hashing.

The difficulty is entirely in choosing the committed bytes, and it is not a matter of taste. A specification directory holds material of several kinds. Programs that decide acceptance. Fixtures they read. Pinned outputs they compare against. Declared resource limits. And, in a system that generates implementations, prose hints addressed to whoever writes the code. Each byte either enters the hash or does not. Both choices can be wrong, and neither failure is visible until something exploits it.

This paper states the conditions precisely, shows that they constrain the committed set from opposite directions, and exhibits violations of each from a working system. Our contributions:

\begin{itemize}\itemsep1pt\parskip0pt
\item Two conditions on the committed byte set --- \emph{determination} and \emph{admissibility} --- with the observation that they force the committed set to \emph{equal} the acceptance procedure's read set, and with a witness for each violation (§\ref{sec:identity}).
\item A third condition on the \emph{reconstruction} procedure, which can otherwise answer a question the specification did not ask (§\ref{sec:preserve}).
\item A composition rule for dependent specifications, and three working forgeries showing what its clauses buy (§\ref{sec:compose}).
\item A necessity argument against purely extensional criteria, with a reconstruction that realized the degenerate member, and the obligation form that excludes it (§\ref{sec:extensional}).
\item A technique for committing \emph{complexity} obligations to a portable identity via growth ratios rather than absolute bounds, with two defects it caught that all functional criteria accepted (§\ref{sec:resource}).
\item Validation at the limit: the system computing these identities reconstructed itself across three generations and two producers under one unchanged identity (§\ref{sec:selfhost}).
\end{itemize}

We state results as \emph{conditions with witnesses}, not theorems with proofs. The acceptance procedure here is an executing program, so its read set is an operational quantity rather than a syntactic one. We demonstrate a violation by exhibiting an input on which the system misbehaves, then showing the misbehavior absent after repair. The lineage is Thompson's trust argument~\cite{thompson}, not formal verification.

\section{The identity construction and its two conditions}
\label{sec:identity}
\subsection{Setup}
A \emph{specification} $S$ is a finite map from paths to byte strings, with each path assigned a declared role. An \emph{implementation} $I$ is a separate finite map, nominally the material a producer must supply. An \emph{acceptance procedure} $A$ evaluates them:
\begin{equation}
A(S,I;e,b,z)\in\{\mathrm{pass},\mathrm{fail},\mathrm{unknown}\},
\end{equation}
with $e$ the execution environment, $b$ a resource envelope, and $z$ the seed of any randomized check. The third outcome separates \emph{missing evidence} from \emph{rejection}: an exceeded declared timeout is a failure, an unavailable platform is not.

An \emph{identity} is a function of a declared subset of the specification's bytes,
\begin{equation}
\rho(S)=H\!\left(\operatorname{canon}\big(\{(p,H(S[p])):p\in\Pi(S)\}\big)\right),
\end{equation}
where $H$ is a cryptographic hash, $\Pi(S)$ is the \emph{committed path set}, and $\operatorname{canon}$ is a canonical encoding. We write $\mathrm{Com}(S)=\{(p,S[p]):p\in\Pi(S)\}$ for the committed content. Implementation bytes are excluded by construction, which is the point: the name must survive replacing them.

Define the \emph{read set} $\mathrm{Rd}(A,S)$: the specification paths whose contents can change the outcome of $A$ on some implementation. Operationally, these are the paths $A$ opens, executes, or receives as arguments, plus those its subprocesses do.

\subsection{Condition 1: determination}
The name is supposed to answer ``would this be judged the same way?'' That requires:

\begin{condition}[Determination]
\label{cond:det}
For all specifications $S_1,S_2$ and all implementations $I$,
\[\rho(S_1)=\rho(S_2)\ \Longrightarrow\ A(S_1,I)=A(S_2,I).\]
\end{condition}

\noindent Determination holds if and only if every byte that can affect the verdict is committed:
\begin{equation}
\mathrm{Rd}(A,S)\subseteq\Pi(S)\quad\text{for every }S.
\label{eq:det}
\end{equation}
The ``only if'' direction is immediate. Suppose some $p\in\mathrm{Rd}(A,S)\setminus\Pi(S)$ influences an outcome. Then two specifications differing only at $p$ share an identity and disagree on a verdict. That is a direct counterexample.

\paragraph{Witness.} The condition failed in our system, in a form worth reporting: the obvious repair is the wrong one. Producer hints are prose attached to a generated output, telling whoever writes it what is wanted. We had removed them from the committed set, for reasons given in §\ref{sec:adm}. But the judging workspace was still built from the \emph{raw} specification, hints included. So an acceptance program that read its own hints could accept differently under one identity. Our documentation asserted this was impossible. A reproduction showed that as implemented it was not (\art{research/audits/finding-a-repro}{counterexample}).

The repair did not add the hints back to $\Pi$. It \emph{shrank the read set}. The judging workspace is now built from the committed content itself --- exactly the bytes the identity serializes. No acceptance program can read what the identity does not cover. A criterion then enforces the property: two specifications differing only in hints share an identity, and must share a verdict \emph{even when the acceptance program goes looking}.

The reason the obvious repair is wrong is the second condition.

\subsection{Condition 2: admissibility}
\label{sec:adm}
An identity that over-commits refuses implementations that satisfy every requirement:

\begin{condition}[Admissibility]
\label{cond:adm}
If $S_1$ and $S_2$ differ only in material that cannot affect any verdict, then $\rho(S_1)=\rho(S_2)$.
\end{condition}

\noindent Equivalently, $\Pi(S)$ must hold nothing outside the specification's semantic content: no presentation, no prose, no incidental ordering. The failure mode here is not a wrong verdict. It is a wrong \emph{name}. An independent implementation satisfies every criterion, produces a different identity, and is rejected as a different thing.

\paragraph{Witness 1 (prose).} Hints were committed at first. Then an independently written packager, conforming in every behavioral respect, phrased its generated hints differently. It minted a different identity for the entire dependency chain. Nothing about acceptance had changed --- only wording. Hints left the committed set, and from that point any conforming packager mints the same names.

\paragraph{Witness 2 (order).} With prose gone, the residue was \emph{sequence}. The step list was committed as a list. Two conforming packagers enumerated the same steps in different orders, one lexicographic and one following glob expansion. They minted different identities, drifting 13 of 20 chain names. The committed encoding now sorts steps by canonical form, making the list a set. Order can no more reject a realization than wording can.

\subsection{The committed set equals the read set}
Conditions~\ref{cond:det} and~\ref{cond:adm} press from opposite sides, and together they are nearly rigid:
\begin{equation}
\mathrm{Rd}(A,S)\ \subseteq\ \Pi(S)\ \subseteq\ \mathrm{Sem}(S).
\label{eq:sandwich}
\end{equation}
Since any byte an acceptance program consults is thereby semantic, $\mathrm{Rd}\subseteq\mathrm{Sem}$ always, and the designer's freedom lies only in the gap. The gap is where the useful material sits: anything the identity should cover because it \emph{might} be consulted later.

The witness above illustrates the practical consequence. Sometimes a byte must be excluded for admissibility, yet the acceptance program can read it. The only sound repair is then to make it \emph{unreadable} --- to separate the context in which an implementation is constructed from the context in which it is judged. In a system that generates implementations, that separation cannot be a documentation convention. It has to be enforced: \emph{the judge sees the committed bytes and nothing else}.

\section{The reconstruction procedure must preserve the question}
\label{sec:preserve}
A third condition concerns the machinery that produces implementations rather than the identity itself. Let $R_T$ be a reconstruction procedure: it materializes a permitted view of $S$, invokes a producer $P$, and seals the result.

\begin{condition}[Question preservation]
If $R_T(S,P)$ reports success with implementation $I$, then $A(S,I)=\mathrm{pass}$ --- for the \emph{supplied} $S$, not for some $S'$ produced during reconstruction.
\end{condition}

\noindent The condition is not automatic, and its violation is easy to miss, because the result is internally consistent. Our reconstruction originally sealed whatever the destination workspace held after the producer ran. Suppose the producer edited a pinned fixture, or the recipe itself. The result verified, its criteria passed, its identity was well-formed --- and it was not the claim requested. The procedure reported success.

The repair snapshots the committed content after materialization, then refuses before sealing if the producer disturbed it. The caller's own deliberate substitutions stay permitted. A meddling producer is now refused in band rather than silently answered.

This is the same structural error as a compiler that accepts a modified source file and reports that the original compiled.

\section{Composition}
\label{sec:compose}
\subsection{The rule}
Specifications depend on other specifications. Let $C$ declare edges $E$, where $(d,o,i)\in E$ means component $d$ supplies its output $o$ to $C$'s input $i$, and let $C$ ship bytes $\beta_C(i)$ at that input. The composed verdict must be
\begin{equation}
A^\ast(C)=A(C)\ \wedge\!\!\bigwedge_{(d,o,i)\in E}\! A\big(D_d,\ \beta_C(i)\ \text{as}\ o\big).
\label{eq:compose}
\end{equation}
Each component's \emph{own} acceptance procedure is re-run on the bytes the dependent actually ships for it. Gates compose; verdicts never carry. A stored ``pass'' from the component's history is not evidence about the dependent's bytes.

Equation~\ref{eq:compose} is easy to state. In our system it was wrong in three places at once. Each is worth exhibiting, because each shows one clause earning its place.

\subsection{Forgery 1: unre-earned data edges}
Edges come in two kinds. A \emph{code} edge supplies a generated output the dependent ships as source; the dependent's recipe names the component. A \emph{data} edge arises when a dependent pins an input whose content matches a component's output --- the dependency is recorded, but no generated step carries it.

Code edges were re-earned. Data edges were content-matched once, when the dependent was sealed, and never again. The attribution then rode in the dependent's metadata --- where the signing construction folded it into chain identity.

The forgery runs in three steps. Seal a component whose criterion produces a data file. Seal a dependent that pins that file and records the attribution. Then overwrite the dependent's copy with different bytes and re-seal.

The dependent's own criterion is blind to this by construction: it never examines provenance. So the claim is internally consistent while asserting an origin the bytes do not have. We ran it against every available implementation, the reference one included. The forged claim passed identity verification, shallow acceptance, \emph{composed} acceptance, and the dependency view. Four verbs, no refusal.

The repair covers every edge no generated step carries. For each, it compares the dependent's shipped input against the resolved component's re-earned output, and refuses naming the attribution. Re-run, the same witness is now rejected at the composed audit while the dependent's own criterion stays green. That is the correct division of labor: only the composed verdict can see this.

\subsection{Forgery 2: an endorsement that does not compose}
A signing ceremony attests that a human reviewed a claim and authorizes it. Ours signed the \emph{chain} identity --- folding every declared component into the signed object --- while re-earning only $A(C)$, the claim's own criteria.

The forgery: corrupt the component code a composed claim ships, so the component's own criterion would reject it while the dependent's criterion cannot see the difference. Then $A(C)=\mathrm{pass}$ and $A^\ast(C)=\mathrm{fail}$ --- and the ceremony signs. The strongest verdict in the system carried the weakest evidence across the composed boundary. The signature named a chain its own composed audit rejects.

The repair re-earns $A^\ast$ whenever a claim declares components, refusing with the broken layers named. For a leaf claim the two audits coincide, so nothing changes. We considered a weaker repair --- recording the \emph{depth} the signature reached --- and rejected it. A reader who must notice a depth field to avoid being misled is not meaningfully protected.

\subsection{Forgery 3: a verb that does not compose}
The first two forgeries concern library-level procedures. The third concerns the interface people actually use. Our command-line audit verb resolved, by default, to composed acceptance --- correct behavior, but behavior only, with no criterion requiring it. A reimplementation was free to wire the verb to the shallow primitive and pass every test we had.

This is not hypothetical. The system is routinely reconstructed by independent producers (§\ref{sec:selfhost}), and nothing in the criteria distinguished the two wirings. The repair drives the dispatch itself against the standard witness. The default must refuse a forged composed claim; the explicitly requested shallow form must accept it. An opt-down the user spells out is honest. A default that silently fails to compose is Forgery 2 again, one layer out.

\subsection{What the three have in common}
Each violation is the \emph{unexercised half of a symmetric contract}. Code edges were re-earned; data edges were not. The signing ceremony's shallow refusal was tested; its composed refusal was not. The verb's denial behavior was pinned; its permission behavior was not.

In every case the criteria exercised one side of a two-sided obligation, leaving an implementation --- or an adversary --- free on the other. This is the most transferable observation in the paper. To audit a specification for omissions, enumerate its contracts' symmetries, then check whether both halves are exercised.

\section{Finite example sets cannot constrain behavior}
\label{sec:extensional}
\subsection{The degenerate member}
Let $S_{\mathrm{ext}}$ be a specification whose criteria consist of a finite set of input/output observations $\{(x_j,y_j)\}_{j\le n}$, compared exactly.

\begin{claimx}[Extensional insufficiency]
The class admitted by $S_{\mathrm{ext}}$ contains an implementation $T$ of size $O(n)$ that stores the observations, answers $y_j$ on each $x_j$, and is unconstrained on every other input.
\end{claimx}

\noindent As mathematics this is trivial. We report it because the degenerate member is not hypothetical. It is what a competent producer builds when the criteria permit it and nothing else is demanded.

We froze a corpus of 465 observations for an expression evaluator and asked for a blind reconstruction. It passed all 465. It was a 465-entry table, with careful input normalization and an error for everything else. Its own header states the position plainly: \emph{``The supplied specification is an extensional set of observations\ldots{} This is not a general interpreter: unspecified program/input combinations raise an error.''} (\art{research/external/jq/evidence/sample1-the-lookup-table.py}{source}). On probes drawn from outside the corpus it matched the reference on 587 of 1{,}395 and refused the remaining 808.

The producer was not evading the criteria. It satisfied them at minimum cost, and documented what it had done. The criteria were the defect.

\subsection{Obligations that a table cannot satisfy}
The repair must rule out the degenerate member, not merely discourage it. Add obligations evaluated over inputs \emph{generated when the acceptance procedure runs}, from entropy drawn at that moment:

\begin{itemize}\itemsep1pt\parskip0pt
\item \emph{Computed relations}: the checker derives the expected result independently for each generated input (indexing, length, concatenation, ordering).
\item \emph{Metamorphic relations}: relations between executions that need no oracle --- composition against sequencing, idempotence, round trips.
\item \emph{Refusal obligations}: generated ill-typed inputs must raise, guarded by a check that valid inputs still pass, so refusing everything cannot satisfy them.
\end{itemize}

A stored table cannot satisfy these. The inputs do not exist when the table is written. This is a necessity argument rather than a recommendation: no finite enlargement of an exact-comparison corpus excludes $T$, and generated obligations do.

Two practical notes. First, the seed must be fresh for each evaluation. A fixed seed turns the obligations into a larger corpus, which §\ref{sec:extensional} then applies to. Second, each obligation has to be checked against the reference before it is committed. Three of our drafted refusal obligations were false. The reference's refusals turned out to be \emph{value}-dependent rather than type-dependent: it reverses the number $0$ to the empty list but raises on $5$, because its definition routes through a length that is zero there. Committed as written, those obligations would have refused conforming implementations.

\subsection{Evidence that the repair steers}
Adding obligations narrows the admitted class by construction. Whether independent reconstructions actually move is a separate, empirical question.

We measured it per behavior, paired before and after. Of the behaviors that diverged from the reference consistently before the obligations were added, 15 diverged in no sample afterwards. One went the other way. That is an exact one-tailed binomial $p\approx2.6\times10^{-4}$, and each closure traces by name to the obligation responsible (\art{research/external/jq/rate-metric-findings.md}{analysis}).

The pairing matters, because the obvious alternative misleads. An earlier aggregate count of divergences per sample suggested a 25\% reduction. A control at an \emph{unchanged} specification returned the same aggregate: between-sample variance exceeded the effect (\art{research/external/jq/generation-4-control-CORRECTION.md}{control}). An aggregate discards the pairing between a behavior before and after, and the pairing is where the signal sits. We treat even the paired result as exploratory --- correlated probes, one producer family, adaptive analysis --- and none of §\ref{sec:identity}--§\ref{sec:resource} depends on it.

\section{Committing complexity obligations}
\label{sec:resource}
\subsection{Why absolute bounds cannot be committed}
Functional criteria accept implementations with arbitrarily bad asymptotics. The natural remedy is to commit a bound: ``this operation completes within $T$ seconds.'' That destroys portability. $T$ is a statement about a machine, so an identity containing it names a specification only some hosts can judge. And a bound loose enough for every host constrains nothing on a fast one.

\subsection{Ratios cancel the host}
Commit a \emph{growth ratio} instead. For an operation parameterized by scale $K$, measured twice on the same host in the same run,
\begin{equation}
\frac{t(2K)}{t(K)}\ \le\ c.
\label{eq:ratio}
\end{equation}
The host's constant factor divides out. What remains constrains the algorithm, not the machine. A slow linear implementation on slow hardware passes; a fast exponential one fails everywhere. The quantity is a property of the code, which is what an identity may legitimately commit. And the measurement is self-calibrating, since both terms come from the same execution.

\subsection{Two defects it caught}
\paragraph{Time.} One reconstruction passed every functional criterion, then exceeded the declared envelope on the full system. Its dependency traversal re-visited each ancestor once per path reaching it, rather than resolving each once. The doubling ratio was $\approx\!95$, against $2.1$--$3.5$ for conforming implementations; a $2^K$ model predicted the interpolated point to within one percent. Committing \eqref{eq:ratio} with $c=10$ moves that refusal from a timeout late in a whole-system audit to a named failure at the component responsible.

\paragraph{Space.} The same defect class on the other axis. A copy procedure duplicated shared ancestors once per reaching path, so a dependency store grew exponentially in chain depth --- in practice to 23\,GB for a 20-layer chain. The committed obligation is again a bound rather than a layout: each dependency's bytes appear $O(1)$ times, and the copy must stand alone once its source is removed. One defect class, two observables; the obligation form is the same.

\subsection{Scope}
A ratio obligation constrains the exponent only at the scales measured. An implementation well behaved to $2K$ and pathological beyond it still passes. A declared absolute envelope therefore remains useful as a backstop --- as an \emph{evaluation condition}, not a committed byte. Obligations of this kind also need calibrating against conforming implementations before they are committed. Our first drafts refused behavior that was legitimately value-dependent, and we abandoned one candidate outright: its refusal condition could not be expressed as a property of types.

\section{Validation at the limit}
\label{sec:selfhost}
The conditions above are claims about a construction that names software by its criteria. The strongest available test is to apply the construction to the system that implements it, and ask whether the identity survives replacing everything it names.

\subsection{The experiment}
A 21-component decomposition of the system was reconstructed in dependency order, each component's producer receiving already-reconstructed dependencies as context and never the original component. Assembled implementations were then judged against the unchanged whole-system identity, with all 37 generated modules substituted.

At one frozen identity, three reconstructions qualified: one from each of two independent producer families, and a complete second generation whose every component was written by one family's producer while every identity computation, sandbox, and reconstruction step was executed by the \emph{other} family's reconstructed code. That third result transfers the operational role, not merely the implementation: the tool performing the work is itself a reconstruction. Its own recursive step then reconstructed a component through its interface.

Measured text similarity between independent realizations is $\approx\!0.12$, and 32 of 37 modules have a distinct syntactic shape in every realization. The identity is therefore not accidentally stable: it survived three complete replacements of the bytes it names.

\subsection{Verdicts as claims}
A verdict about a claim is itself evidence, and evidence stored as a log is unverifiable. Each qualifying reconstruction is sealed as a claim whose \emph{own} acceptance procedure re-executes the checks --- identity under substitution, held-out tests, operational steps --- so that the verdict is re-earned rather than asserted (\art{research/harness/closure/trial_claim.py}{builder}). A claim cannot be sealed unless its criteria pass, so only a qualifying trial seals at all; the construction refused several, including one that caught a dangling path reference in our own committed criteria.

The regress this invites terminates deliberately. Judges are auditable by the machinery they judge with, and the one operation never automated is \emph{authorization}: adopting a new obligation, and thus moving an identity, requires a human signature. Every obligation discussed in this paper carries one, with the refusal that selected it recorded alongside (\art{research/metrics/genome-map.md}{register}).

\subsection{What is retained}
Self-hosting claims require an explicit partition. Reconstructed: the identity computation, sandboxing, reconstruction, transport, acceptance execution, and packaging. Retained: the language runtime, the operating system, the producer, the adapter invoking it, and the experiment driver choosing component order. Source exclusion is a property of the producer's workspace, which is testable; it is not a claim about the producer's training data, which is not.

One condition of our own stopping rule is open: no operator outside the authors' environment has repeated the judgment. Continuous integration re-executes the criteria on six platform configurations we do not administer, but it runs our workflow, and the distinction matters. A protocol and signed-receipt format for an external operator are specified in advance (\art{research/harness/closure/outside.md}{protocol}); until one is returned, the self-hosting result is a laboratory result, and we mark it so.

\section{Related work}
\paragraph{Trust in self-reproduction.} Thompson showed that a tool's capacity to reproduce itself says nothing about its trustworthiness~\cite{thompson}. Diverse double-compiling tests source-executable correspondence using an independent compiler under stated assumptions~\cite{wheeler}. Our question is adjacent but distinct: not whether an executable corresponds to its source, but whether fresh source can be constructed from committed requirements and then operate as the constructing tool. We eliminate neither the trusted runtime nor the producer.

\paragraph{Reproducible and bootstrappable builds.} Reproducible builds recover identical artifacts from identical inputs~\cite{builds}; bootstrappable builds minimize the trusted binary seed~\cite{bootstrappable}. Both fix the implementation and vary nothing. We fix the \emph{requirements} and allow the implementation to vary, which is why the conditions of §\ref{sec:identity} arise at all: a reproducible build needs no admissibility condition, because it admits exactly one artifact.

\paragraph{Verified compilation.} CakeML bootstraps a verified compiler inside a theorem prover, relating an implementation to a formal semantics~\cite{cakeml}. That provides guarantees of a kind a finite executable criterion cannot. Our setting is the one where the specification is executable and incomplete, and the contribution is the soundness of the \emph{naming and composition} machinery under that weakness.

\paragraph{Specification mining.} Dynamic invariant inference proposes properties a program appears to maintain~\cite{daikon}. It generalizes from executions of an existing implementation; we generate new implementations and treat their disagreements as evidence about the specification. The two compose: inferred invariants are candidate obligations, and a reconstruction violating one indicates either a missing obligation or an overgeneralized inference.

\paragraph{Generated-input testing.} Property-based and metamorphic testing supply the obligation forms of §\ref{sec:extensional}~\cite{quickcheck,metamorphic}, and differential testing the disagreement-finding method~\cite{csmith}. Our use is unusual in placing such obligations \emph{inside a committed identity}, which is what forces the fresh-seed requirement: an obligation whose inputs are fixed is a corpus, and §\ref{sec:extensional} applies to it.

\paragraph{Multiversion programming.} The idea of accepting any of several independent implementations is old~\cite{avizienis}, as is the finding that independent implementations fail alike~\cite{knight}. The latter recurred here: six reconstructions of a deliberately underspecified parser, three from each of two producer families, all passed its criteria and all broke the same consumer in the same way (\art{research/harness/substitution}{harness}). What has changed is cost --- generating a version is now cheap enough to run the experiment repeatedly --- which was the practical obstacle, not the conceptual one.

\section{Limitations}
The conditions of §\ref{sec:identity} are stated over an operational read set. Determining $\mathrm{Rd}(A,S)$ exactly requires knowing what an arbitrary program reads; we approximate it by construction (materializing only committed bytes into the judging context) rather than by analysis, which is sound for the resulting system but does not yield a decision procedure for an arbitrary one.

The forgeries of §\ref{sec:compose} demonstrate that specific clauses of \eqref{eq:compose} are necessary. They do not establish that \eqref{eq:compose} is sufficient, and we know of no argument that it is; further edge kinds may require further clauses, as the data-edge case did.

The empirical material comes from one system, whose criteria were authored by the method's designer and revised repeatedly in response to the behavior of the two producer families used to probe them. Specification and probe population co-evolved. That is the right way to develop a boundary and a confound when measuring one; a specification authored independently of the producers would be a stronger test. The external pilot of §\ref{sec:extensional} is partial and its target is widely represented in producer training data, so its agreement results are uninformative and only its divergences are used.

Finally, §\ref{sec:selfhost}'s open condition is the one a reader should weigh most heavily. Three generations is not an induction, and nothing in our machinery references generation depth; a chain that held three times may fail the fourth, and no external operator has yet repeated the judgment.

\section{Conclusion}
Hashing requirements instead of bytes gives software a name that survives reimplementation. The construction is trivial; its soundness is not. The committed byte set must contain everything the acceptance procedure can read, or one name admits contradictory verdicts, and nothing the specification does not mean, or conforming implementations are refused for presentation. Those bounds pin the committed set to the read set, and when the two conflict the repair is to restrict what the judge can see rather than to widen what the name covers.

Composition needs its own rule, and ours was wrong in three places that forged claims exploited --- each the unexercised half of a two-sided obligation. Finite example sets cannot constrain behavior outside their support, and a reconstruction demonstrated the degenerate member rather than merely permitting it; obligations over judgment-time inputs are the necessary repair. Complexity obligations can enter a portable identity as growth ratios but not as absolute bounds, and in that form they caught an exponential that every functional criterion accepted.

That these conditions hold in practice was tested at the limit, by the system that computes the identities: across three generations and two independent producers, with 0.12 measured similarity between realizations, one name held while everything it names was replaced.

\balance
\section*{Artifact availability}
All witnesses, criteria, generated implementations, and measurement records are in the public repository at snapshot \texttt{5318acaac34e\-f99eeba9ff06\-4a995a6313cc70ec}~\cite{repo}; links throughout resolve at that commit. The forgery witnesses of §\ref{sec:compose} are executable, as is the counterexample of §\ref{sec:identity}. Reconstructed implementations are destroyed after measurement by design; one early probe record was overwritten by a later run reusing its path and was regenerated exactly from the preserved generated source, the probe inputs being seed-determined (587 agreements and 808 refusals over 1{,}395 probes, matching the figures recorded at the time). No other number here depends on a file that was overwritten, and no new producer runs or external replay were performed for this manuscript.

\begin{thebibliography}{12}\small
\bibitem{repo} R. Zamora-Resendiz. \emph{Reticuli}. Repository snapshot, October 2026. \url{https://github.com/rz4/reticuli/tree/5318acaac34ef99eeba9ff064a995a6313cc70ec}.
\bibitem{thompson} K. Thompson. Reflections on trusting trust. \emph{Communications of the ACM}, 27(8):761--763, 1984. doi:10.1145/358198.358210.
\bibitem{wheeler} D. A. Wheeler. \emph{Fully Countering Trusting Trust through Diverse Double-Compiling}. Ph.D. dissertation, George Mason University, 2009. \url{https://arxiv.org/abs/1004.5534}.
\bibitem{cakeml} R. Kumar, M. O. Myreen, M. Norrish, and S. Owens. CakeML: A verified implementation of ML. \emph{POPL}, 2014. doi:10.1145/2535838.2535841.
\bibitem{daikon} M. D. Ernst, J. H. Perkins, P. J. Guo, S. McCamant, C. Pacheco, M. S. Tschantz, and C. Xiao. The Daikon system for dynamic detection of likely invariants. \emph{Science of Computer Programming}, 69(1--3):35--45, 2007. doi:10.1016/j.scico.2007.01.015.
\bibitem{quickcheck} K. Claessen and J. Hughes. QuickCheck: A lightweight tool for random testing of Haskell programs. \emph{ICFP}, pp. 268--279, 2000. doi:10.1145/351240.351266.
\bibitem{metamorphic} T. Y. Chen, S. C. Cheung, and S. M. Yiu. \emph{Metamorphic Testing: A New Approach for Generating Next Test Cases}. Technical Report HKUST-CS98-01, 1998. \url{https://arxiv.org/abs/2002.12543}.
\bibitem{csmith} X. Yang, Y. Chen, E. Eide, and J. Regehr. Finding and understanding bugs in C compilers. \emph{PLDI}, 2011. doi:10.1145/1993498.1993532.
\bibitem{avizienis} A. Avi\v{z}ienis. The N-version approach to fault-tolerant software. \emph{IEEE Transactions on Software Engineering}, SE-11(12):1491--1501, 1985. doi:10.1109/TSE.1985.231893.
\bibitem{knight} J. C. Knight and N. G. Leveson. An experimental evaluation of the assumption of independence in multiversion programming. \emph{IEEE Transactions on Software Engineering}, SE-12(1):96--109, 1986. doi:10.1109/TSE.1986.6312924.
\bibitem{builds} C. Lamb and S. Zacchiroli. Reproducible builds: Increasing the integrity of software supply chains. 2021. \url{https://arxiv.org/abs/2104.06020}.
\bibitem{bootstrappable} J. Nieuwenhuizen, L. Courtès, et al. Bootstrappable builds and the full-source bootstrap. \url{https://bootstrappable.org/}, accessed October 2026.
\end{thebibliography}
\end{document}
